%% ma: L12-B12-0001
%% lop: 12
%% chuong: 4
%% bai: 12
%% dang: 12.12.1
%% loai: TLN
%% muc: VD
%% dap_an: "43"
%% nguon: "Tổng hợp VDC THPT 2026-2027 (PDF giáo viên gửi), Câu 2"
%% kien_thuc: "Bài 12 (tích phân), Bài 2 (GTNN, chương 1 — chương trước) — không vượt chương"
%% ghi_chu: "Hình gốc có ký hiệu góc vuông thừa tại O, đã bỏ khi vẽ lại"
%% trang_thai: chinh-thuc
\begin{ex}
Trên đoạn $[0;8]$ cho hàm số
\[f(x)=\begin{cases}-x(x-4), & 0\le x<4,\\ x-4, & 4\le x\le 8.\end{cases}\]
Với số thực $a$ thỏa mãn $0\le a\le 4$, giá trị nhỏ nhất của $\displaystyle\int_a^{a+4}f(x)\,\mathrm dx$ là $\dfrac qp$, trong đó $p,q$ là hai số tự nhiên nguyên tố cùng nhau. Tính $p+q$.
\begin{center}
\begin{tikzpicture}[scale=.6]
  \draw[truc] (-.6,0)--(9,0) node[below]{$x$};
  \draw[truc] (0,-.6)--(0,5) node[left]{$y$};
  \draw[phu] (2,0)--(2,4) (0,4)--(8,4) (8,0)--(8,4);
  \draw[thick,domain=0:4,samples=50] plot (\x,{-\x*(\x-4)});
  \draw[thick] (4,0)--(8,4) node[pos=.55,above left]{$y=f(x)$};
  \path (0,0) node[below left]{$O$} (2,0) node[below]{$2$} (4,0) node[below]{$4$} (8,0) node[below]{$8$} (0,4) node[left]{$4$};
\end{tikzpicture}
\end{center}
\shortans{43}
\loigiai{Với $0\le a\le4$: $F(a)=\displaystyle\int_a^4(4x-x^2)\,\mathrm dx+\int_4^{a+4}(x-4)\,\mathrm dx=\frac{a^3}{3}-\frac{3a^2}{2}+\frac{32}{3}$.
$F'(a)=a^2-3a=0\Leftrightarrow a=0$ hoặc $a=3$. So sánh $F(0)=\frac{32}{3}$, $F(3)=\frac{37}{6}$, $F(4)=8$ được $\min F=\frac{37}{6}$, nên $p=6$, $q=37$, $p+q=43$.}
\end{ex}
